\documentclass[10pt]{article}

\usepackage[utf8]{inputenc}

\usepackage[T1]{fontenc}

\usepackage{amsmath}

\usepackage{amsfonts}

\usepackage{amssymb}

\usepackage[version=4]{mhchem}

\usepackage{stmaryrd}

\usepackage{graphicx}

\usepackage[export]{adjustbox}

\graphicspath{ {./images/} }



\begin{document}

Однако, во-первых, при достаточно большом $n$ ни одну из функций әтого класса, несмотря на то, что он содержит почти всо булевы функции, не удается явно указать, а во-вторых, этот класс состоит из наиболее сложно реализуемых функций, к-рые, по-видимому, не встречаются в задачах прикладного характера. Поәтому такие функции не представляют большого интереса. И снова появляется необходимость в видоизменении постановки С. 3.



При этом следует иметь в виду, что функции, для к-рых С. 3. актуальна, составляют ничтожную долю всех булевых функций, и что любой әффективный универсальный метод синтеза (если гипотеза о необходимости перебора верна в полном объеме) для нек-рых из них заведомо даст далеко не самые простые схемы. Описать множество таких функций невозможно, но проблема классификации булевых функций для целей синтеза возникает. Такой подход к постановке С. 3. связан с выделением важных классов функций, созданием для них специальных методов синтеза, каждый из к-рых ориентирован на функции определенного класса, а также с оценкой качества получающихся при этом схем.



Эфффективно решать такие задачи позволяет п р и нци п лока льно го коди р о в ани я. Этот принцип предполагает такое кодирование функций из рассматриваемого класса, при к-ром значение функции можно вычислять по сравнительно небольшой части кода. Если при әтом удается достаточно просто реализовать нек-рые вспомогательные операторы, то принцип локального кодирования дает асимптотически наилучший метод синтеза для рассматриваемого



т а б л и ц 2 класса, причем в доволь-



\begin{center}

\begin{tabular}{c|c}

\hline

$\begin{array}{c}\text { Число точек, } \\ \text { в к-рых } \\ \text { рункция } \\ \text { равна 1 }\end{array}$ & Т а б л и ц а 2 \\

\hline

$\left(\log _{2} n\right)^{c}, c>1$ & $\sim n\left(\log _{2} n\right)^{c-1}$ \\

\hline

$n^{c}, c>0$ & $\sim \frac{n^{c+1}}{(c+1) \log _{2} n}$ \\

\hline

$2^{n^{c}, 0<c<1}$ & $\sim n^{1-c 2^{n} c}$ \\

\hline

$2^{c n}, 0<c<1$ & $\sim \frac{1-c}{c} 2^{c n}$ \\

\hline

$\frac{2^{n}}{n^{c}}, c>0$ & $\sim c \frac{2^{n} \log _{2} n}{n^{c+1}}$ \\

\hline

\end{tabular}

\end{center}



но широких пределах справедливо соотношение



$L\left(N_{n}\right) \sim \rho \frac{\log _{2}\left|N_{n}\right|}{\log _{2} \log _{2}\left|N_{n}\right|}$, где $N_{n}-$ подмножество тех функций из рассматриваемого класса, к-рые зависят от заданных $n$ переменных, а



\[

L\left(N_{n}\right)=\max L(f),

\]



где максимум берется по всем булевым функциям $f$ из $N_{n}$. (здесь $\rho=1$ ).



Ковкретный вид, который может иметь это асимптотическое соотношение, удобнее всего



проиллюстрировать на классах булевых функций, принимающих значение 1 в заданном числе точек. Наиболее характерные примеры приведены в таблице 2.



На основе принципа локального кодирования были получены хорошие методы синтеза и для весьма узких классов функций (напр., для симметрич. Функций), т. е. для тех функций, к-рые с практической точки зрения наиболее интересны. Однако в данном случае нижняя оценка, получающаяся мощностным методом, оказывается слишком слабой, для того чтобы по ней судить о качестве схем. По мере сужения класса рассматриваемых функций в какой-то момент это неизбежно должно было произойти, поскольку предельным случаем узкого класса является класс, содержащий не более чем по одной функции от каждого числа переменных. Фактически это означает, что рассматривагтся отдельные конкретные функции.



Здесь возникают две проблемы: 1) построить для заданной конкретной функции $f$ схему с возможно меньшим числом элементов и тем самым получить верхнюю оценку для $L(f) ; 2)$ получить возможно лучІпую нижнюю оценку для $L(f)$. Среди этих задач первая обычно бывает легче, так как для ее решения достаточно, в конечном счете, построить одну схему, а вторая требует в каком-то смысле просмотра всех схем, реализующих функцию $f$. По этой причине возникло целое направление, связанное с поиском методов получения нижних оценок для конкретных функций.



Получение нижних оценок, линейных относительно $n$ - числа аргументов функции (или в общем случае относительно объема входной информации), обычно все-таки не вызывает больших трудностей. Получение более сильных нижних оценок - значительно более сложная задача. Первыми примерами в этом направлении были нижняя оценка порядка $n^{1,5}$ для формул в базисе $\{\&, \vee,-\}$, реализующих линейную функцию (сумму по модулю 2) $n$ аргументов, и нижняя оценка порядка $n^{2} / \log _{2}^{2} n$ для нормальных алгорифмов Маркова, к-рые обращают слова длины $n$. Наиболее высокие нижние оценки, к-рые к настоящему времени (1984) удалось получить, имеют порядок $n^{2}$, если не считать оценок, получающихся при очень сильных ограничениях на класс у. с. (напр., неполнота системы әлементов) либо при использовании столь мощных средств для задания функции, что они фактически содержат в себе перебор всех функций.



$\mathrm{O}$ нетривиальности исследования сложности реализации конкретных функций говорит и то, что многие задачи, такие, как умножение целых чисел, умножение матриц, распознавание симметрии слова на многоленточных машинах Тьюринга, для к-рых ожидались сравнительно высокие нижние оценки, оказались на самом деле проще, чем предполагалось раньше.



При минимизации других параметров схем - задержек, вероятности сбоя и т. д.- возникают шримерно такие же задачи. При этом удается установить определенную связь между различными параметрами, причем благодаря моделированию одних объектов на других такую связь часто удается установить между различными параметрами у. с. из разных классов. Таким образом, полученные в С. 3. результаты не являются изолированными, а тесно переплетаются друг с другом и часто переносятся из одной области в другую.



Многое из сказанного выше относится и к бесконечным моделям у. с. Однако с постановкой задачи синтеза, а тем более с ее решением дело обстоит значительно сложнее.



Лum.: [1] JI у п а н о в О. Б., «Проблемы кибернетики», 1965 , в. 14, с. $31-110 ;$ [2] Я б л о н с к и й С. В., там же,

1959 , в. 2 , $75-121$.



\begin{center}

\includegraphics[max width=\textwidth]{2023_07_08_928353487eda67c712b7g-1}

\end{center}



\[

y=\sin x

\]



Область определения - вся числовая прямая, область значений - отрезок $[-1 ; 1]$; С.- функция нечетная, периодическая с периодом $2 \pi$. С. и косинус связаны формулой



\[

\sin ^{2} x+\cos ^{2} x=1

\]



С. и косеканс связаны формулої



\[

\sin x=\frac{1}{\operatorname{cosec} x}

\]



Производная С.:



\[

(\sin x)^{\prime}=\cos x

\]



Интеграл от С.:



\[

\int \sin x d x=-\cos x+C

\]



С. разлагается в степенной ряд:



\[

\sin x=x-\frac{x^{3}}{3 !}+\frac{x^{5}}{5 !}-\ldots,-\infty<x<\infty .

\]





\end{document}